\documentclass[UTF8]{ctexart}
\usepackage{geometry,booktabs,longtable,array,tabularx}
\geometry{a4paper,margin=2.0cm}
\setlength{\emergencystretch}{3em}
\setlength{\tabcolsep}{3pt}
\renewcommand{\arraystretch}{1.12}
\title{JiuwenSwarm v4 可审计科研论文半流程报告}
\author{JiuwenSwarm/UCR}
\date{}
\begin{document}
\maketitle
\begin{abstract}
本文汇总三个主题的科研半流程演示。系统检索公开资料元数据，执行 JiuwenSwarm/UCR 本地证据实验，生成研究计划和报告，并保存 Claim--Evidence 账本。本文不把流程基准包装成领域科学结论；正式引用、实验真实性、研究结论和最终提交资格仍需人工确认。
\end{abstract}
\section{统一边界}
UCR 的标签是 \texttt{process\_evidence\_only}，它表示系统的流程证据和审计能力，不等于上下文工程、记忆引擎或自我进化的领域效果。所有来源默认处于待审核状态。当前 PDF 是由本机 XeLaTeX 编译生成的正式查看产物，源文件和实验记录同时保留在提交包中。
\section{矩阵概览}
\begin{longtable}{@{}>{\raggedright\arraybackslash}p{0.30\textwidth}>{\raggedright\arraybackslash}p{0.42\textwidth}>{\raggedleft\arraybackslash}p{0.14\textwidth}@{}}
\toprule
主题 & 状态 & Provider 调用\\\\
\midrule
\endfirsthead
\toprule
主题 & 状态 & Provider 调用\\\\
\midrule
\endhead
context-\allowbreak{}engineering & \texttt{\scriptsize completed\_\allowbreak{}pending\_\allowbreak{}human\_\allowbreak{}review} & 3 \\
memory-\allowbreak{}engine & \texttt{\scriptsize completed\_\allowbreak{}pending\_\allowbreak{}human\_\allowbreak{}review} & 3 \\
self-\allowbreak{}evolution & \texttt{\scriptsize completed\_\allowbreak{}pending\_\allowbreak{}human\_\allowbreak{}review} & 3 \\
\bottomrule
\end{longtable}
\section{Agent 上下文工程}
\textbf{摘要：}本报告检验程序级证据护栏能否提高 Agent 执行报告的证据一致性。三臂实验均完成且 return\_\allowbreak{}code=0。no-\allowbreak{}rail 的 UCR=0.5556（5 / 9），prompt-\allowbreak{}only 与 full-\allowbreak{}rail 的 UCR 均为 0.0（0 / 4）。full-\allowbreak{}rail 的 rail.enabled=true、attempts=1、accepted=true。三臂 task\_\allowbreak{}success\_\allowbreak{}rate 均为 1.0。CFR 与 NFR 状态为 not\_\allowbreak{}measured\_\allowbreak{}in\_\allowbreak{}ucr\_\allowbreak{}scenario，不纳入结论。\\
\textbf{结论：}在本实验条件下，full-\allowbreak{}rail 的 UCR 不高于 no-\allowbreak{}rail 与 prompt-\allowbreak{}only，与假设方向一致；但 prompt-\allowbreak{}only 与 full-\allowbreak{}rail 的 UCR 相同，无法区分提示词与程序级护栏的增量效应。该结论仅基于过程证据，需人工确认并补充重复运行。\\
\textbf{限制：}UCR 仅为过程证据，不能推断任务质量或真实完成度。CFR 与 NFR 状态为 not\_\allowbreak{}measured\_\allowbreak{}in\_\allowbreak{}ucr\_\allowbreak{}scenario，不得纳入结论。citation\_\allowbreak{}allowed=false 的材料不得作为正式引用。prompt-\allowbreak{}only 与 full-\allowbreak{}rail 的 UCR 相同，单次运行不足以支持护栏增量效应，需重复运行或更大样本。
\section{Agent 记忆引擎}
\textbf{摘要：}本报告检验结构化记忆记录是否有助于 Agent 保持审计上下文一致。三个 arm（no-\allowbreak{}rail、prompt-\allowbreak{}only、full-\allowbreak{}rail）均完成且 return\_\allowbreak{}code=0，任务成功率均为 1.0。UCR 仅作为 process\_\allowbreak{}evidence\_\allowbreak{}only：no-\allowbreak{}rail 为 0.5556 (5 / 9)，prompt-\allowbreak{}only 与 full-\allowbreak{}rail 均为 0.0 (0 / 4)。CFR / NFR 在 UCR 场景中未测量。材料最高 relevance\_\allowbreak{}score 为 0.32 且 citation\_\allowbreak{}allowed=false，故不作正式引用。当前证据不足以支持或拒绝领域效果假设。\\
\textbf{结论：}三个 arm 均完成且任务成功率均为 1.0，但 UCR 仅反映流程证据：no-\allowbreak{}rail 为 0.5556 (5 / 9)，prompt-\allowbreak{}only 与 full-\allowbreak{}rail 均为 0.0 (0 / 4)。当前证据不足以支持或拒绝“结构化记忆记录有助于 Agent 保持审计上下文一致”的领域效果假设，需人工确认 UCR 定位、citation\_\allowbreak{}allowed=false 材料处理、CFR / NFR 后续测量及结构化记忆记录的操作化定义。\\
\textbf{限制：}UCR 仅为 process\_\allowbreak{}evidence\_\allowbreak{}only，不能解释为领域效果。CFR / NFR 在 UCR 场景中未测量。材料最高 relevance\_\allowbreak{}score 为 0.32 且 citation\_\allowbreak{}allowed=false，citation\_\allowbreak{}coverage=0.0，不能作为正式引用。claim-\allowbreak{}002、claim-\allowbreak{}003、claim-\allowbreak{}005、claim-\allowbreak{}007、claim-\allowbreak{}009、claim-\allowbreak{}018、claim-\allowbreak{}019、claim-\allowbreak{}020 未支持，需人工复核。allow\_\allowbreak{}paper=false，human\_\allowbreak{}intervention\_\allowbreak{}required=true。
\section{Agent 自我进化}
\textbf{摘要：}本报告检验带证据的反思循环能否避免把未执行任务写成已完成。三臂实验（no-\allowbreak{}rail、prompt-\allowbreak{}only、full-\allowbreak{}rail）均完成且返回码为0。no-\allowbreak{}rail的UCR=0.5556（5 / 9），prompt-\allowbreak{}only与full-\allowbreak{}rail的UCR均为0.0（0 / 4）。任务成功率三臂均为1.0。CFR与NFR在UCR场景中未测量。所有来源relevance\_\allowbreak{}score低且citation\_\allowbreak{}allowed=false，citation\_\allowbreak{}coverage=0.0，因此不进入正式引用。结论仅限过程证据，不证明真实自我进化。\\
\textbf{结论：}在当前基准下，prompt-\allowbreak{}only与full-\allowbreak{}rail的UCR均为0.0，no-\allowbreak{}rail为0.5556，提示证据检查可能改善审计表达，减少未执行任务被写成已完成的比例。但该结果仅为过程证据，不证明真实自我进化。citation\_\allowbreak{}allowed=false的来源不进入正式引用。需补充CFR / NFR测量并人工确认后方可支持更强结论。\\
\textbf{限制：}UCR仅作为过程证据指标报告，不解释为真实自我进化能力。CFR与NFR在UCR场景中未测量，无法支持更强结论。citation\_\allowbreak{}coverage=0.0，所有来源relevance\_\allowbreak{}score低且citation\_\allowbreak{}allowed=false，不能支持正式引用（claim-\allowbreak{}018至claim-\allowbreak{}020为背景待审）。no-\allowbreak{}rail中5个未执行操作的审计状态为pending\_\allowbreak{}human\_\allowbreak{}review，需人工确认。三臂任务集与操作定义是否一致需人工确认。
\section{编译与审核说明}
该统一报告和三个主题报告都需要人工阅读。编译命令为：
\begin{verbatim}
xelatex -interaction=nonstopmode -halt-on-error paper.tex
xelatex -interaction=nonstopmode -halt-on-error paper.tex
\end{verbatim}
最终状态保持为 \texttt{prepared\_not\_uploaded}。
\end{document}
